\section{Template Method}
\label{sec:pat-template-method}

\problemstatement{Define the skeleton of an algorithm in a base-class
method, deferring some steps to subclasses, so subclasses can redefine
certain steps without changing the algorithm's overall structure.}

\begin{patternbox}[When You Reach for It]
The trigger is \emph{``the same overall procedure with a few steps that
differ''} -- data importers that all parse/validate/save but read
different formats, report generators, test set-up/tear-down. When the
outline is fixed but a few steps vary, fix the outline and vary the steps.
\end{patternbox}

\subsection*{Understanding the pattern}

A base class provides a non-virtual \texttt{templateMethod()} that calls a
sequence of steps in a fixed order. Some steps are concrete (shared), some
are \texttt{virtual}/abstract ``hooks'' that subclasses override. The
invariant sequence lives in one place; only the variable steps are
subclassed. This inverts control: the base class calls down into the
subclass (``don't call us, we'll call you'').

\begin{ideabox}[Fix the Skeleton, Override the Steps]
Keep the algorithm's order in a base method and expose just the varying
steps as overridable hooks. Subclasses cannot accidentally reorder or skip
the algorithm -- they only fill in the blanks.
\end{ideabox}

\subsection*{C++ implementation}

\begin{lstlisting}
#include <bits/stdc++.h>
using namespace std;

class DataProcessor {
public:
    void process() {                            // template method (fixed order)
        read();
        auto data = parse();
        cout << "Processing " << data << " rows\n";
        save();
    }
    virtual ~DataProcessor() = default;
protected:
    virtual void read() = 0;                     // varying steps
    virtual int  parse() = 0;
    void save() { cout << "Saved to DB\n"; }     // shared step
};

class CsvProcessor : public DataProcessor {
protected:
    void read() override { cout << "Read CSV file\n"; }
    int  parse() override { return 42; }
};

int main() {
    CsvProcessor p;
    p.process();                                // skeleton runs, steps vary
    return 0;
}
\end{lstlisting}

\begin{complexitybox}[Trade-offs]
\textbf{Pros.} Removes duplicated algorithm structure; enforces a fixed
sequence; subclasses supply only what differs. \textbf{Cons.} Relies on
inheritance (less flexible than composition); the base--subclass contract
about which hooks to override can be subtle; changing the skeleton affects
all subclasses.
\end{complexitybox}

\begin{takeawaybox}
Template Method locks an algorithm's structure in a base method and lets
subclasses override only the variable steps -- inheritance-based inversion
of control. Its composition-based cousin is Strategy: Template Method
varies steps via subclassing, Strategy varies the whole algorithm via an
injected object.
\end{takeawaybox}
